최근 자연어 처리(Natural Language Processing) 분야에서 대형 언어 모델(Large Language Models, LLMs)은 눈부신 발전을 이룩하며 인간의 언어 능력을 모방하는 데 성공했습니다. 텍스트 요약, 번역, 코드 작성 등 다양한 영역에서 상용화 수준의 성과를 보이고 있으나, 고도의 논리적 엄밀성과 다단계 계산이 요구되는 수학적 추론(Mathematical Reasoning) 영역에서는 여전히 치명적인 한계를 드러내고 있습니다. 모델이 텍스트를 생성할 때 발생하는 `환각(Hallucination)' 현상과 산술적 연산 능력의 부재는 복잡한 수식과 증명을 요구하는 문제에서 오답을 정당화하는 확증 편향(Confirmation Bias)으로 이어집니다.

본 연구는 이러한 LLM의 근본적인 한계를 극복하기 위해, AI가 인간 수학자의 사고 과정을 단계별로 모방할 수 있도록 설계된 \textbf{OMI (Orchestrated Math Interpreter)} 시스템을 제안합니다. 이 시스템은 AI가 머릿속으로 모든 것을 계산하게 두는 대신, 파이썬(Python) 코드와 같은 외부 도구를 적극적으로 활용할 수 있는 환경을 제공합니다. 

본 보고서에서는 시스템의 뼈대가 되는 \textbf{5단계 심층 추론 파이프라인(5-Stage Deep Reasoning Pipeline)}을 제안합니다. 이 파이프라인은 문제의 도메인을 분류하는 Labeling 단계부터, 필요한 수학적 지식을 불러오는 Retrieval, 문제 스케일에 따라 시뮬레이션과 이론적 해법을 결정하는 Router, 파이썬 샌드박스를 통해 실제 수식을 계산하는 Execution, 그리고 결과를 다각도로 검증하는 Verification 단계로 구성됩니다. 

또한, 단일 모델의 한계를 돌파하기 위해 다수결 투표(Majority Voting)와 다중 에이전트 토론(Multi-Agent Debate)이라는 앙상블 기법을 도입하여 수학적/물리적 무결성을 획기적으로 향상시켰습니다. 나아가 AI 모델의 추론 능력을 고도화하기 위한 학습 데이터의 오염 문제를 해결하고자, 정답으로부터 문제를 파생시키는 \textbf{역공학(Reverse Engineering) 기반의 데이터 생성 기법(Open-Reverse-Math)}을 최초로 고안 및 적용하였습니다.

최종적으로, 이 모든 복잡한 백엔드 아키텍처와 AI의 사고 과정을 일반 사용자와 학습자가 직관적으로 이해할 수 있도록 시각화한 `투명한 유리 상자(Glass-box)' 형태의 웹 대시보드를 구축하였습니다. 이를 통해 OMI 프로젝트가 단순한 질의응답 시스템을 넘어, 향후 논리와 추론이 필요한 모든 산업 및 교육 도메인으로 확장 가능한 차세대 Agentic AI의 표준 아키텍처임을 입증하고자 합니다.
